% Source: https://artofproblemsolving.com/community/c3h1885879_a_fast_way_to_check_large_numers_multiple
% Stable identity: aops-topic-1885879-q01
(a) Since $10\equiv1\pmod9$, every power of $10$ is $\equiv1\pmod9$, so every positive integer is congruent to the sum of its digits modulo $9$. Repeating, $r(m)\equiv m\pmod9$. Moreover $r(m)$ is one of $1,2,\dots,9$, because a positive integer has positive digit sum. These nine values are pairwise incongruent modulo $9$, so $r(m)$ is \emph{the} number in $\{1,\dots,9\}$ congruent to $m$ modulo $9$; in particular $r(m)=r(n)$ whenever $m\equiv n\pmod9$. Now
\[
ab\equiv r(a)\,r(b)\pmod9,
\]
hence $r(ab)=r\bigl(r(a)\,r(b)\bigr)$. $\blacksquare$

(b) $r(325)=r(10)=1$ and $r(675)=r(18)=9$, so $r\bigl(r(325)r(675)\bigr)=r(9)=9$. Also $r(219375)=r(27)=9$. They agree.

(c) No. The test only checks the product modulo $9$. A wrong result with the same remainder modulo $9$ passes, for example $219735$ (two digits swapped) also has digital root $9$. Agreement is a check, not a proof; disagreement does prove an error.
